% Formulario TEEA 4MAT — Argomento 1: segnali elettrici continui e alternati
% Per rigenerare il PDF (dalla cartella strumenti, due volte per il numero di pagine):
%   pdflatex -output-directory=../pdf -jobname=01-formulario-segnali-continui-alternati formulario.tex
% ---------------------------------------------------------------------------
%  Stile comune dei formulari TEEA (pdfLaTeX, solo pacchetti standard
%  presenti in TeX Live e MiKTeX). Testo in Computer Modern Sans.
% ---------------------------------------------------------------------------
\documentclass[11pt,a4paper]{article}
\usepackage[a4paper, top=14mm, bottom=16mm, left=14mm, right=14mm, headheight=0pt, footskip=8mm]{geometry}
\usepackage[utf8]{inputenc}
\usepackage{amsmath, amssymb}
\renewcommand{\familydefault}{\sfdefault}
\newcommand{\ohm}{\ensuremath{\Omega}}
\newcommand{\uA}{\ensuremath{\mu}}
\newcommand{\per}{\ensuremath{\cdot}}
\usepackage{xcolor, tikz, multicol, booktabs, array, fancyhdr, lastpage, tabularx}
\usepackage[most]{tcolorbox}
\usetikzlibrary{arrows.meta}

\definecolor{verde}{HTML}{159957}
\definecolor{blu}{HTML}{155799}
\definecolor{inchiostro}{HTML}{1B2936}
\definecolor{tenue}{HTML}{606C71}
\definecolor{fondo}{HTML}{EEF3F0}
\definecolor{tensione}{HTML}{2152C0}
\definecolor{corrente}{HTML}{B3205F}
\definecolor{resistenza}{HTML}{0B7A62}
\color{inchiostro}

\setlength{\parindent}{0pt}
\setlength{\columnsep}{7mm}
\renewcommand{\arraystretch}{1.25}

% unità di misura in tondo e virgola decimale
\newcommand{\U}[1]{\,\mathrm{#1}}
\newcommand{\virg}{{,}}

% riquadro di un argomento
\newtcolorbox{blocco}[1]{enhanced, breakable=false, colback=white, colframe=verde, boxrule=0.6pt,
  leftrule=3pt, arc=1.5pt, left=6pt, right=6pt, top=4pt, bottom=5pt,
  title={\bfseries #1}, coltitle=verde, colbacktitle=fondo, fonttitle=\normalsize,
  toptitle=1pt, bottomtitle=1pt, before skip=0pt, after skip=4.5mm}

% formula grande centrata
\newcommand{\F}[1]{\par\smallskip{\centering\Large$\displaystyle #1$\par}\smallskip}
% legenda dei simboli
\newcommand{\leg}[1]{{\small\color{tenue}#1}\par}

% intestazione del formulario
\newcommand{\intestazione}[3]{%
  {\color{verde}\bfseries\small #1}\par\vspace{1mm}
  {\fontsize{19}{23}\selectfont\bfseries #2}\par\vspace{1mm}
  {\small\color{tenue} #3}\par\vspace{2mm}
  {\color{verde}\rule{\linewidth}{1.6pt}}\par\vspace{4mm}}

\pagestyle{fancy}
\fancyhf{}
\renewcommand{\headrulewidth}{0pt}
\fancyfoot[L]{\scriptsize\color{tenue}Prof. Pasquale Naclerio \per{} TEEA 4MAT \per{} Formulario}
\fancyfoot[R]{\scriptsize\color{tenue}Pagina \thepage\ di \pageref{LastPage}}

\begin{document}

\intestazione{TEEA --- 4MAT \per{} Argomento 1 \per{} FORMULARIO}
             {Segnali elettrici continui e alternati}
             {Prof. Pasquale Naclerio \per{} Manutenzione e Assistenza Tecnica}

\begin{multicols}{2}

\begin{blocco}{Sinusoide}
\F{v(t) = V_M \cdot \sin(\omega t + \varphi)}
\F{i(t) = I_M \cdot \sin(\omega t + \varphi)}
\begin{tabularx}{\linewidth}{@{}c X l@{}}
\toprule
$V_M,\ I_M$ & ampiezza (valore di picco) & V, A\\
$V_{pp}$ & valore picco-picco & V\\
$T$ & periodo & s\\
$f$ & frequenza & Hz\\
$\omega$ & pulsazione & rad/s\\
$\varphi$ & fase iniziale & rad, $^\circ$\\
\bottomrule
\end{tabularx}
\end{blocco}

\begin{blocco}{Periodo, frequenza, pulsazione}
\F{f = \frac{1}{T} \qquad T = \frac{1}{f}}
\F{\omega = 2\pi f = \frac{2\pi}{T} \qquad f = \frac{\omega}{2\pi}}
\end{blocco}

\begin{blocco}{Valori della sinusoide}
\F{V_{pp} = 2\,V_M \qquad V_M = \frac{V_{pp}}{2}}
\F{V_{\mathrm{eff}} = \frac{V_M}{\sqrt{2}} \qquad V_M = \sqrt{2}\,V_{\mathrm{eff}}}
\end{blocco}

\begin{blocco}{Gradi e radianti}
\F{\alpha_{\mathrm{rad}} = \alpha^\circ \cdot \frac{\pi}{180} \qquad \alpha^\circ = \alpha_{\mathrm{rad}} \cdot \frac{180}{\pi}}
\smallskip
{\footnotesize\setlength{\tabcolsep}{3pt}
\begin{tabularx}{\linewidth}{@{}l *{8}{>{\centering\arraybackslash}X}@{}}
\toprule
gradi & $0^\circ$ & $30^\circ$ & $45^\circ$ & $60^\circ$ & $90^\circ$ & $180^\circ$ & $270^\circ$ & $360^\circ$\\
rad & 0 & $\frac{\pi}{6}$ & $\frac{\pi}{4}$ & $\frac{\pi}{3}$ & $\frac{\pi}{2}$ & $\pi$ & $\frac{3\pi}{2}$ & $2\pi$\\
\bottomrule
\end{tabularx}}
\end{blocco}

\columnbreak

\begin{blocco}{Fasore}
\F{\bar{V} = V_M \angle \varphi}
\leg{lunghezza $= V_M$ \quad angolo iniziale $= \varphi$ \quad velocità di rotazione $= \omega$ (senso antiorario)}
\end{blocco}

\begin{blocco}{Sfasamento}
\F{\Delta\varphi = \varphi_1 - \varphi_2}
\F{\Delta\varphi = \frac{\Delta t}{T} \cdot 360^\circ \qquad \Delta\varphi = \frac{\Delta t}{T} \cdot 2\pi}
\F{\Delta t = \frac{\Delta\varphi}{\omega}}
\leg{$\Delta\varphi > 0$: in anticipo \quad $\Delta\varphi < 0$: in ritardo \quad $0$: in fase}
\leg{$180^\circ$: opposizione di fase \quad $90^\circ$: quadratura}
\end{blocco}

\begin{blocco}{Lettura all'oscilloscopio}
\F{n = \text{div intere} + \text{parti} \cdot 0\virg2}
\F{V = n_{\mathrm{vert}} \cdot \mathrm{V/div}}
\F{t = n_{\mathrm{oriz}} \cdot \mathrm{Time/div}}
\leg{1 divisione = 5 parti \quad 1 parte = 0,2 div}
\end{blocco}

\begin{blocco}{Rete elettrica}
\F{f = 50\U{Hz} \qquad T = 20\U{ms}}
\F{\omega \approx 314\U{rad/s}}
\F{V_{\mathrm{eff}} = 230\U{V} \quad V_M \approx 325\U{V}}
\end{blocco}

\begin{blocco}{Multipli e sottomultipli}
\begin{tabularx}{\linewidth}{@{}X X@{}}
$1\U{kHz} = 10^{3}\U{Hz}$ & $1\U{ms} = 10^{-3}\U{s}$\\
$1\U{MHz} = 10^{6}\U{Hz}$ & $1\,\mu\mathrm{s} = 10^{-6}\U{s}$\\
\end{tabularx}
\end{blocco}

\end{multicols}

\end{document}
